\section{Conclusion}
We have developed a complete, modular, and explainable multi-stage deep learning pipeline for brain tumor segmentation and classification from 2D MRI slices. The ResUNet segmentation backbone achieves localization performance (Dice 0.7431, IoU 0.6558) and effectively focuses downstream classification while enabling robust interpretability. 

The dual-input ResNet50 classifier combines segmentation-guided ROI crops with full-image context and attains 99.00\% accuracy and 98.95\% F1-score. This helps address the accuracy-interpretability trade-off by recovering global anatomical context lost in pure ROI cropping, while still preserving explicit tumor localization.

Grad-CAM visualizations and a quantitative stage-consistency IoU metric verify that the classifier focuses on the same tumor regions identified by segmentation, providing an objective safeguard against spurious correlations. For textual explanations, MedGemma outperforms other VLMs (BLIP-2, InstructBLIP, CLIP) in clinically relevant terminology and composite summary quality; the hybrid MedGemma VLM + template approach reduces the risk of prediction-inconsistent summaries and keeps textual explanations aligned with model predictions.

\paragraph{Limitations and Future Work.}
Several limitations remain. First, the current pipeline processes 2D MRI slices independently and does not model volumetric context across adjacent slices. Future work could extend the segmentation stage to 3D U-Net or transformer-based architectures, which may improve localization of small or irregular tumors. Second, the VLM summary module still relies on a fallback template when generation is malformed or inconsistent with the classifier prediction. Incorporating calibrated uncertainty estimation or selective prediction could help the system decide when to generate free text and when to fall back to a safer template.

Third, our evaluation is based on publicly available datasets with relatively high-quality annotations, so additional validation on external datasets and other tumor types is needed to assess generalizability. The same multi-stage explainable framework could also be adapted to other medical imaging tasks, such as lung nodule assessment, diabetic retinopathy grading, or histopathology analysis. Finally, Grad-CAM and the segmentation-consistency metric provide useful evidence, but they do not guarantee clinical correctness. Future user studies with radiologists are needed to evaluate whether the combined visual and textual explanations improve diagnostic trust and decision making. Reducing computational overhead through model distillation or a shared backbone would also make the system more practical for deployment.